\documentclass[11pt]{article}

\usepackage[margin=1in]{geometry}
\usepackage{amsmath, amssymb}
\usepackage{graphicx}
\usepackage{listings}
\usepackage{xcolor}
\usepackage{hyperref}
\usepackage{booktabs}

\lstset{
    basicstyle=\ttfamily\small,
    breaklines=true,
    frame=single,
    numbers=left,
    numberstyle=\tiny,
    keywordstyle=\color{blue},
    commentstyle=\color{gray},
    tabsize=4
}

\title{High Performance Parallel Computing \\ Homework Template}
\author{Your Name(s)}
\date{\today}

\begin{document}

\maketitle

\section{Problem Statement}
Briefly restate the problem being solved.

\section{Approach}
Describe the parallelization strategy (e.g. domain decomposition, threading
model, communication pattern) and any relevant design decisions.

\section{Implementation}
Highlight the key parts of the code. Reference \texttt{src/main.c} rather
than pasting the whole file; use \texttt{listings} for short excerpts.

\begin{lstlisting}[language=C, caption={Example excerpt}]
int main(int argc, char *argv[]) {
    printf("hello\n");
    return 0;
}
\end{lstlisting}

\section{Results}
Present timing/scaling data, plots, or tables here.

% \begin{figure}[h]
%     \centering
%     \includegraphics[width=0.7\textwidth]{figures/speedup.pdf}
%     \caption{Speedup vs. number of processes/threads.}
% \end{figure}

\section{Discussion}
Interpret the results: does performance match expectations? What are the
bottlenecks?

\section{Conclusion}
Summarize findings.

\end{document}
